\documentclass[11pt,a4paper]{article}

% ===================== PAQUETES =====================
\usepackage[utf8]{inputenc}
\usepackage[spanish]{babel}
\usepackage[T1]{fontenc}
\usepackage[margin=2.5cm]{geometry}
\usepackage{graphicx}
\usepackage{float}
\usepackage{hyperref}
\usepackage{booktabs}
\usepackage{amsmath}
\usepackage{fancyhdr}

\graphicspath{{capturas/}}

% ===================== CONFIG HYPERREF =====================
\hypersetup{
    colorlinks=true,
    linkcolor=blue,
    urlcolor=blue,
    pdftitle={Etapa 0: simulación previa en Python -- Informe de presentación},
    pdfauthor={Lina Robsana Jamioy Tisoy, Tomás Rocha Castaño}
}

% ===================== ENCABEZADOS =====================
\pagestyle{fancy}
\fancyhf{}
\fancyhead[L]{Robótica -- Proyecto final}
\fancyhead[R]{Etapa 0: simulación en Python}
\fancyfoot[C]{\thepage}

% ===================== COMANDO PARA FIGURAS =====================
% Uso: \figura{nombre_del_archivo}{descripción}{etiqueta}
\newcommand{\figura}[3]{%
    \begin{figure}[H]
        \centering
        \includegraphics[width=0.85\linewidth]{#1}
        \caption{#2}
        \label{#3}
    \end{figure}
}

% ======================================================
\begin{document}

% ===================== PORTADA =====================
\begin{titlepage}
    \centering
    \vspace*{2cm}

    {\Large \textbf{UNIVERSIDAD NACIONAL DE COLOMBIA}}\\[0.3cm]
    {\large Sede Manizales}\\[2cm]

    {\Huge \textbf{Plataforma robótica de laberinto con control autónomo mediante ROS 2 y visión artificial}}\\[0.5cm]
    {\Large Informe de presentación: Etapa 0}\\[0.3cm]
    {\large Simulación previa en Python}\\[2cm]

    \rule{\linewidth}{0.5mm}\\[0.4cm]
    {\large \textbf{Curso:} Robótica}\\
    {\large \textbf{Código:} MAN-4101136-1-2026-02}\\[0.4cm]
    \rule{\linewidth}{0.5mm}\\[1.5cm]

    {\Large \textbf{Integrantes:}}\\[0.4cm]
    {\large Lina Robsana Jamioy Tisoy}\\
    {\large Tomás Rocha Castaño}\\[2cm]

    {\large \textbf{Profesor:}}\\
    {\large Jaime Enrique Arango Castro}\\[2cm]

    \vfill
    {\large \today}

\end{titlepage}

\tableofcontents
\newpage

% ===================== 1. INTRODUCCIÓN =====================
\section{Introducción}
El proyecto consiste en guiar una esfera a través de un laberinto de $7 \times 7$ celdas, inclinando una plataforma de dos grados de libertad. El sistema combina percepción visual, planificación de trayectorias y control, coordinados con ROS 2 y simulados en Gazebo.

Antes de construir el modelo tridimensional, decidimos validar la lógica del sistema en una simulación bidimensional escrita en Python. Esta primera etapa permite detectar errores de planificación, de control o de modelado físico sin que se mezclen con problemas propios de Gazebo o de ROS 2. Cada módulo se probó de forma independiente, con pruebas automáticas y con gráficas que documentan su comportamiento.

% ===================== 2. OBJETIVOS =====================
\section{Objetivos de la Etapa 0}
\begin{enumerate}
    \item Verificar que la representación lógica del laberinto coincide con el diseño del plan.
    \item Implementar y probar el planificador A* sobre esa representación.
    \item Validar el modelo físico de la esfera rodando sobre la plataforma inclinada.
    \item Diseñar y probar un controlador PD que siga la ruta sin salirse del corredor, y evaluar su robustez ante el ruido de medición.
\end{enumerate}

% ===================== 3. QUÉ SE VA A REALIZAR =====================
\section{Qué se va a realizar}
La etapa se divide en cinco módulos de Python, cada uno con sus pruebas automáticas y sus gráficas:

\begin{itemize}
    \item \textbf{\texttt{maze.py}:} define la matriz del laberinto, las celdas libres y las reglas de vecindad.
    \item \textbf{\texttt{planner.py}:} calcula la ruta de inicio a salida con A* y compara dos criterios de desempate.
    \item \textbf{\texttt{physics.py}:} modela la aceleración de la esfera según la inclinación de la plataforma.
    \item \textbf{\texttt{controller.py}:} convierte la posición y la ruta en inclinaciones de los ejes X y Y, con un controlador PD.
    \item \textbf{\texttt{simulate.py}:} ejecuta el recorrido completo, guarda los datos en CSV y genera una animación.
\end{itemize}

Además, se realiza una prueba de robustez: el controlador se ejecuta 30 veces por cada nivel de ruido de medición, para comprobar que sigue llegando a la meta.

% ===================== 4. RESULTADOS ESPERADOS =====================
\section{Resultados esperados}
Antes de ejecutar las pruebas, se definieron los siguientes resultados esperados:

\begin{itemize}
    \item La matriz del laberinto tiene inicio en la celda $(0,0)$ y salida en la celda $(6,6)$.
    \item A* encuentra una ruta de 12 pasos que no atraviesa paredes.
    \item Con una inclinación de $5^\circ$, la aceleración de la esfera es de $0{,}61$ m/s$^2$, y tarda unos $0{,}36$ s en recorrer una celda de 40 mm desde el reposo.
    \item El controlador llega a la meta con un sobreimpulso menor al $10\,\%$, sin colisiones con las paredes.
    \item Con un ruido de medición de 1 mm, el recorrido se completa en todas las corridas.
\end{itemize}

% ===================== 5. DESARROLLO Y RESULTADOS =====================
\section{Desarrollo y resultados}

\subsection{Mapa lógico del laberinto}
La primera verificación consiste en comprobar que el programa imprime la misma matriz definida en el plan. En la salida de texto, \texttt{S} marca el inicio, \texttt{G} la salida, \texttt{\#} las paredes y \texttt{.} las celdas libres.

\figura{etapa0_01_maze_terminal}
{Salida de \texttt{python src/maze.py}. La matriz coincide con la definida en el plan, con inicio en $(0,0)$ y salida en $(6,6)$.}
{fig:maze_terminal}

La figura~\ref{fig:maze_terminal} muestra la matriz tal como la imprime el programa. Esta es la verificación visual de que la representación lógica coincide con el diseño, y sirve como punto de partida para todos los módulos siguientes.

\figura{etapa0_03_pytest_ok}
{Resumen de las pruebas del módulo \texttt{maze}: seis pruebas pasan.}
{fig:pytest_maze}

La figura~\ref{fig:pytest_maze} muestra las seis pruebas del módulo \texttt{maze}: inicio y salida, celdas libres y paredes, límites del tablero y vecinos de $S$ y $G$. Durante el desarrollo, tres de estas pruebas fallaron por dos causas: la función devolvía un tipo de dato de NumPy en lugar de un booleano nativo, y una expectativa de la prueba no coincidía con la matriz. Ambas se corrigieron antes de continuar.

\figura{etapa0_04_laberinto_grafica}
{Representación gráfica del laberinto: inicio en verde, salida en rojo y paredes en gris oscuro.}
{fig:maze_grafica}

La figura~\ref{fig:maze_grafica} presenta el mismo laberinto de forma gráfica. Facilita la lectura de las rutas posibles y es la base de las figuras de planificación.

\subsection{Planificación con A*}
El planificador compara dos criterios de desempate cuando dos rutas tienen el mismo costo: \texttt{primera}, que toma el nodo insertado primero, y \texttt{menor\_h}, que toma el nodo más cercano a la meta.

\figura{etapa0_05_planner_salida}
{Salida de \texttt{python src/planner.py}: ambos criterios encuentran una ruta de 12 pasos y expanden 12 nodos.}
{fig:planner_salida}

La figura~\ref{fig:planner_salida} muestra que, en este laberinto, ambos criterios producen la misma ruta. Esto es esperable, porque la ruta más corta no tiene empates y el desempate no cambia el resultado.

\figura{etapa0_06_pytest_planner}
{Resumen de las pruebas de \texttt{maze} y \texttt{planner}: 14 pruebas pasan.}
{fig:pytest_planner}

La figura~\ref{fig:pytest_planner} muestra las catorce pruebas acumuladas hasta este punto. Verifican que la ruta empieza en $S$ y termina en $G$, que no atraviesa paredes, que se reporta correctamente cuando la salida está bloqueada, y que el criterio inválido genera un error.

\figura{etapa0_07_planner_ruta}
{Ruta calculada por A* sobre el laberinto, con 12 pasos.}
{fig:planner_ruta}

La figura~\ref{fig:planner_ruta} muestra la ruta encontrada. Avanza por la fila superior, baja por la columna 2, continúa por la fila 2 hasta la columna 4, desciende a la fila 3 y termina por la columna 6 hasta llegar a $G$.

\figura{etapa0_08_planner_expandidos}
{Nodos expandidos por cada criterio de desempate, en el laberinto y en un mapa abierto de $7 \times 7$.}
{fig:planner_expandidos}

La figura~\ref{fig:planner_expandidos} muestra la diferencia entre los criterios. En el laberinto ambos expanden 12 nodos. En el mapa abierto, \texttt{primera} expande 48 nodos y \texttt{menor\_h} solo 12. El desempate por cercanía a la meta reduce el trabajo del algoritmo.

\figura{etapa0_09_planner_escala}
{Nodos expandidos en mapas abiertos de $5 \times 5$ a $11 \times 11$, para los dos criterios.}
{fig:planner_escala}

La figura~\ref{fig:planner_escala} muestra cómo crece el trabajo del algoritmo al aumentar el tamaño del mapa. Con \texttt{primera}, el número de nodos expandidos pasa de 24 a 120. Con \texttt{menor\_h}, coincide con la longitud de la ruta, que va de 8 a 20 pasos. La ventaja de \texttt{menor\_h} se mantiene al crecer el mapa.

\subsection{Modelo físico de la esfera}
La esfera rueda sin deslizar sobre la plataforma. Con un ángulo $\theta$, su aceleración a lo largo del eje es $a = \frac{5}{7}\, g \sin\theta$. La simulación integra esta ecuación con un paso de 1 ms y se compara con la fórmula teórica.

\figura{etapa0_10_physics_aceleracion}
{Aceleración de la esfera según la inclinación: curva teórica y puntos obtenidos por integración numérica.}
{fig:physics_aceleracion}

La figura~\ref{fig:physics_aceleracion} muestra que los puntos simulados coinciden con la curva teórica. Esto confirma que la integración numérica reproduce el modelo. Con $5^\circ$ la aceleración es de $0{,}61$ m/s$^2$, y con $8^\circ$ es de $0{,}975$ m/s$^2$.

\figura{etapa0_11_physics_caida}
{Posición en $x$ de una esfera que parte del reposo, con inclinaciones de $5^\circ$ y $8^\circ$. Los triángulos marcan el tiempo teórico para recorrer 40 mm.}
{fig:physics_caida}

La figura~\ref{fig:physics_caida} muestra el tiempo que tarda la esfera en recorrer una celda de 40 mm. Los triángulos marcan los valores teóricos: $0{,}362$ s con $5^\circ$ y $0{,}286$ s con $8^\circ$. Las curvas cruzan la línea de 40 mm en esos instantes.

\figura{etapa0_12_pytest_physics}
{Resumen de las pruebas acumuladas hasta el módulo \texttt{physics}: 20 pruebas pasan.}
{fig:pytest_physics}

La figura~\ref{fig:pytest_physics} muestra las veinte pruebas acumuladas. Las de física verifican la aceleración con $5^\circ$, que la esfera no se mueve sin inclinación, el sentido del movimiento según el signo del ángulo, la independencia de los dos ejes y el tiempo de cruce de una celda.

\subsection{Controlador PD}
El controlador calcula la inclinación de cada eje como $\theta = K_p\, e - K_d\, v$, donde $e$ es el error de posición y $v$ la velocidad filtrada. Las ganancias se obtienen a partir de una frecuencia natural $\omega_n$ y un amortiguamiento $\zeta$:
\begin{equation}
    K_p = \frac{\omega_n^2}{k g}, \qquad K_d = \frac{2\zeta\omega_n}{k g}, \qquad k g = \frac{5}{7}\, g.
\end{equation}
Se eligió $\omega_n = 4$ rad/s y $\zeta = 0{,}7$, que dan $K_p \approx 2{,}28$ rad/m y $K_d \approx 0{,}80$ s/m.

\figura{etapa0_13_controller_escalon}
{Respuesta a un escalón de 40 mm para tres frecuencias naturales, con $\zeta = 0{,}7$.}
{fig:controller_escalon}

La figura~\ref{fig:controller_escalon} compara tres frecuencias naturales. Con $\omega_n = 2$ rad/s, la esfera tarda 1,98 s en llegar; con 4 rad/s, 1,03 s; y con 6 rad/s, 0,71 s. El sobreimpulso se mantiene cerca del $3{,}5\,\%$ en los tres casos, por debajo del $10\,\%$ exigido. Se eligió $\omega_n = 4$ rad/s porque es más rápido que 2 rad/s y evita la saturación que aparecería con 6 rad/s. Con un error de 40 mm, $\omega_n = 6$ rad/s pediría una inclinación de unos $11{,}8^\circ$, por encima del límite de $8^\circ$.

\figura{etapa0_14_controller_ruta}
{Trayectoria de la esfera sobre el laberinto, siguiendo los puntos de la ruta calculada por A*.}
{fig:controller_ruta}

La figura~\ref{fig:controller_ruta} muestra la trayectoria de la esfera siguiendo la ruta de A*. La esfera llega a $G$ sin colisiones en 7,72 s. Los puntos naranjas son los nodos de la ruta, y la trayectoria azul los sigue por el camino central del laberinto.

\figura{etapa0_15_controller_comandos}
{Arriba: inclinaciones comandadas en los ejes X y Y, con el límite de $\pm 8^\circ$ marcado. Abajo: error de seguimiento en cada eje.}
{fig:controller_comandos}

La figura~\ref{fig:controller_comandos} permite comprobar dos condiciones. Primero, las inclinaciones se mantienen dentro de $\pm 8^\circ$, así que el controlador no satura. Segundo, el error de seguimiento tiende a cero al final del recorrido, lo que indica que la esfera se estabiliza en la meta.

\figura{etapa0_16_pytest_controller}
{Resumen de las pruebas acumuladas hasta el controlador: 26 pruebas pasan.}
{fig:pytest_controller}

La figura~\ref{fig:pytest_controller} muestra las veintiséis pruebas acumuladas hasta este punto. Incluyen la fórmula de las ganancias, la saturación a $8^\circ$, la detección de colisiones con paredes y con el borde del tablero, la respuesta al escalón y el recorrido completo del laberinto sin colisiones.

\subsection{Recorrido completo}
El módulo \texttt{simulate.py} une el planificador y el controlador en un recorrido completo de S a G. El resultado se guarda en \texttt{datos/etapa0\_recorrido\_nominal.csv}, con el tiempo, la posición, las inclinaciones y el error en cada paso de 1 ms. Además, genera una animación de la esfera en el laberinto.

\figura{etapa0_17_recorrido}
{Cuadro final de la animación del recorrido completo de S a G, con la esfera sobre la trayectoria y el tiempo transcurrido.}
{fig:recorrido}

La figura~\ref{fig:recorrido} muestra el último cuadro de la animación. La esfera llega a la salida sin colisiones y el registro completo queda guardado en el archivo CSV para la validación posterior.

\subsection{Robustez ante ruido}
La prueba de robustez simula la medición de la posición con un ruido gaussiano de desviación $\sigma$. Por cada nivel de ruido se ejecutan 30 corridas con semillas distintas.

Durante la primera ejecución de esta prueba, la corrida número 22 con $\sigma = 5$ mm chocó con una pared a los 2,34 s. La causa probable era el cambio de punto de la ruta: con una tolerancia de 15 mm, la esfera empezaba a girar antes de llegar al centro de la celda, y con su inercia podía invadir una pared. La tolerancia se redujo a 5 mm, y en la nueva ejecución ninguna de las 30 corridas chocó. El costo fue un recorrido más lento: el tiempo nominal pasó de 6,26 s a 7,72 s.

\begin{table}[H]
    \centering
    \begin{tabular}{@{}rccc@{}}
        \toprule
        \textbf{Ruido $\sigma$ [mm]} & \textbf{Éxito} & \textbf{Tiempo medio [s]} & \textbf{Desviación [s]} \\
        \midrule
        0   & 100 \% & 7{,}72 & 0{,}00 \\
        0{,}5 & 100 \% & 7{,}53 & 0{,}02 \\
        1   & 100 \% & 7{,}30 & 0{,}04 \\
        2   & 100 \% & 6{,}90 & 0{,}15 \\
        5   & 100 \% & 6{,}30 & 0{,}33 \\
        \bottomrule
    \end{tabular}
    \caption{Resultados de robustez con 30 corridas por nivel de ruido.}
    \label{tab:robustez}
\end{table}

La tabla~\ref{tab:robustez} resume los resultados. Con ruido de hasta 5 mm, el controlador llega a la meta en todas las corridas. Sin embargo, el tiempo medio disminuye al aumentar el ruido. La causa es que el cambio de punto se decide con la posición medida: con ruido, la esfera a veces cambia de punto antes de llegar al centro de la celda y recorta la esquina. Este efecto acorta el tiempo, pero no provoca colisiones en las pruebas realizadas.

\figura{etapa0_18_robustez}
{Éxito y tiempo de recorrido según el nivel de ruido de medición, con 30 corridas por nivel.}
{fig:robustez}

La figura~\ref{fig:robustez} presenta los mismos resultados en forma gráfica. La barra de éxito se mantiene en el $100\,\%$ en todos los niveles, y la curva de tiempo muestra la reducción descrita arriba. Como la prueba usa 30 corridas por nivel, los resultados deben leerse como una evidencia inicial y no como una estadística definitiva.

\figura{etapa0_19_pytest_robustez}
{Resumen de las pruebas acumuladas hasta la prueba de robustez: 27 pruebas pasan.}
{fig:pytest_robustez}

La figura~\ref{fig:pytest_robustez} muestra las veintisiete pruebas acumuladas. La última verifica que, con ruido de 1 mm, la esfera completa el recorrido sin colisiones.

% ===================== 6. CONCLUSIONES =====================
\section{Conclusiones}
La Etapa 0 cumple los resultados esperados. La matriz del laberinto, el planificador, el modelo físico y el controlador superan sus pruebas automáticas, y el recorrido completo de S a G se realiza sin colisiones. El controlador mantiene el éxito del $100\,\%$ hasta 5 mm de ruido de medición en 30 corridas por nivel.

Quedan dos observaciones para el análisis posterior. La primera es que la tolerancia de cambio de punto influye directamente en las colisiones y en el tiempo de recorrido. La segunda es que el tiempo medio disminuye con el ruido, lo que indica que el controlador es sensible a la forma en que se decide cambiar de punto.

\subsection{Siguiente etapa}
La Etapa 1 construye el modelo de la plataforma, el laberinto y la esfera en Gazebo. Los parámetros de masa, fricción y paso de simulación se ajustarán para que el comportamiento coincida con los resultados de esta etapa.

\end{document}
